
\documentclass[a4paper,12pt]{letter}
\usepackage[utf8]{inputenc}
\usepackage[english,slovene]{babel}
%\usepackage{blindtext}
%\usepackage{pdflscape}
\usepackage[tikz]{bclogo}
\usepackage{qrcode,marvosym}
\usepackage{pgf,tikz,pgfplots}
\pgfplotsset{compat=1.14}
\usepackage{mathrsfs}
\usetikzlibrary{arrows}
\usepackage{graphicx}
\newcommand{\ds}{\displaystyle}
\usepackage{fancybox}
\usepackage{amsfonts}
\usepackage{amsmath}
%\usepackage{color}
\usepackage{wallpaper}
\usepackage{hyperref}
\usepackage[cp1250]{inputenc}
\usepackage{array}%%%%%%%%%%%%%%%%%%%%%%%%%%%%%%%%%%%%%%%%%%%%
\newcolumntype{M}[1]{>{\centering\arraybackslash}m{#1}}%%%%
\newcolumntype{N}{@{}m{0pt}@{}}%%%%%%%%%%%%%%%%%%%%%%%%%%%%
   \usepackage{fancyhdr,enumerate}

 \newcommand{\ora}{\overrightarrow}
 \renewcommand{\headrulewidth}{3pt}
 \renewcommand{\headrule}{\hbox to\headwidth{%
   \color{gray}\leaders\hrule height \headrulewidth\hfill}}
    \renewcommand{\footrulewidth}{2pt}
    \renewcommand{\footrule}{\hbox to\headwidth{%
  \color{brown}\leaders\hrule height \footrulewidth\hfill}}
\usepackage{gensymb}
\usepackage{amssymb} 

  \usepackage{fancyhdr}
  \textwidth 20 true cm
  \oddsidemargin -2 true cm
  \evensidemargin -2 true cm
  \topmargin -3.5 true cm
  \textheight 27.5 true cm
  \linespread{1.2}
    \renewcommand{\headrulewidth}{0.3pt}
   \renewcommand{\footrulewidth}{1.9pt}
  \definecolor{qqwuqq}{rgb}{1,0,0 }
  \definecolor{qqqqff}{rgb}{0,0,1}
  \definecolor{qqffqq}{rgb}{0,1,0}
  \definecolor{ffqqqq}{rgb}{1,0,0}
\definecolor{zzuxux}{rgb}{0,0,0}%{0.92,0.03,0.03}
\usepackage{mdframed}
\usepackage{multicol}
 \definecolor{bgblue}{RGB}{0,0,253}
\usepackage{xcolor}
\usepackage{eurosym}
%%%%%%%%%%%%%%%%%%%%%%%%%%%%%%%%%%%%%%%%%%%%%%%%%%%%%%%%
\newcounter{tocke}
\setcounter{tocke}{0}
\usepackage{tikz-3dplot}
%%%%%%%%%%%%%%%%%%%%%%%%%%%%%%%%%%%%%%%%%%%%%%%%%%%%%%%%%%
\mdfdefinestyle{theoremstyle}{%
linecolor=brown!40,linewidth=1pt,%
frametitlerule=true,%
frametitlebackgroundcolor=brown!50,
innertopmargin=\topskip,
}

\mdfdefinestyle{theoremstyle1}{%
linecolor=brown,middlelinewidth=1pt,%
frametitlerule=true,%
apptotikzsetting={\tikzset{mdfframetitlebackground/.append style={%
shade,left color=black!40, right color=gray!10},
mdfbackground/.append style={
top color=white!100!white,
bottom color=black!5
},
}}, 
frametitlerulecolor=gray,
frametitlerulewidth=2pt,
innertopmargin=\topskip,
innerbottommargin=\topskip,
}
\mdtheorem[style=theoremstyle1]{naloga}{Naloga}
\mdtheorem[style=theoremstyle]{resitev}{Analiza Naloge}
\mdtheorem[style=theoremstyle]{kriterij}{\v{S}tevilo dose\v{z}enih to\v{c}k na testu}
%%%%%%%%%%%%%%%%%%%%%%%%%%%%%%%%%%%%%%%%%%%%%%%%%%%%%%%%%%
\usepackage[table]{xcolor}
\usepackage{titlesec}
\usepackage{venndiagram}
\usepackage[printwatermark]{xwatermark}
\usepackage{xcolor}
\newwatermark[allpages,color=brown!2,angle=45,scale=5,
xpos=-5,ypos=0]{matej.info}


%%%%%%%%%%%%%%%%%%%%%%%%%%%%%%%%%%%%%%%%%%%%%%%%%%%%%%%%%%
\titleformat{\subsection}
  {\normalfont\Large\bfseries\color{brown}}
  {\thesubsection}
  {1em}
  {}

\titleformat{\subsubsection}
  {\normalfont\large\bfseries\color{brown}}
  { \thesubsubsection}
  {1em}
  {}
%%%%%%%%%%%%%%%%%%%%%%%%%%%%%%%%%%%%%%%%%%%%%%%%%%%%%%%%%%
\begin{document}
\pagestyle{fancy}
\renewcommand\headrulewidth{0pt}
\lhead{}\chead{}\rhead{}
\rfoot{\vspace*{-2\baselineskip}}


\renewcommand{\vec}{\overrightarrow}

%%%%%%%%%%%%%%%%%%%%%%55\Large

\everymath{\displaystyle}
\begin{minipage}{20cm}
\begin{bclogo}[couleur=brown!40,
  arrondi=0,ombre=true, couleurOmbre=gray!50,
logo =\bcplume,
  noborder=true]
{\textrm{{\Large{Test 2.1.} \Huge{}:
%
%
%%%%%%%%%%%%%%%%%%%%%%%%%%%%%%%%%%%%%%%%%
%%%%%%%%%%%%%%%%%%%%%%%%%%%%%%%%%%%%%%%%%
% TUKAJ VNESI VSEBINO TESTA
\large{Eksponentna in logaritemska funkcija. Zaporedja.} \hfill $G-3$}}}
%%%%%%%%%%%%%%%%%%%%%%%%%%%%%
\end{bclogo}
\end{minipage}
\begin{minipage}{20cm}
\begin{bclogo}[couleur=brown!40,ombre=true, couleurOmbre=gray!50,
  arrondi=0,
logo =\bcplume,
  noborder=true]
 \setlength{\parskip}{}{}
{\textsc{{\large{Ime in priimek:}
%%%%%%%%%%%%%%%%%%%%%%%%%%%%%%%%%%%%%%%%%
% TUKAJ VNESI VSEBINO TESTA
\vspace{0.2cm} \rule{15cm}{0.2mm}\hfill }}}
%%%%%%%%%%%%%%%%%%%%%%%%%%%%%%%%%%%%%%%%%
\end{bclogo}
\end{minipage}


\def\tocke2{3+3+3}
\begin{naloga}[\hfill \quad $\tocke2$
\qquad \rightsquigarrow \vert a.\qquad|b.\qquad|c.\qquad|] \addtocounter{tocke}{\tocke2}
Reši enačbo:
\vspace{-0.5cm}
\begin{multicols} {3}
   a) $\log_3(2x-1)=2$

b) $3^x=6$

c) $\sqrt{2}=2^{x-\frac{1}{2}}$ 
\end{multicols}

\end{naloga}






\newpage
\def\tocke18{5+1}
\begin{naloga}[\hfill \quad $\tocke18$
\qquad \rightsquigarrow \vert a.\qquad\vert b.\qquad|
c.\qquad|] \addtocounter{tocke}{\tocke18}

a) Nariši graf funkcije $f(x)=3^{x}-3$, tako da najprej izračunaš začetno vrednost in ničlo ter določiš asimptoto.

b) Izračunaj vrednost točke $A(2,y)$, če leži na grafu ter jo označi.
\end{naloga}

\begin{tikzpicture}[line cap=round,line join=round,>=triangle 45,x=1.0cm,y=1.0cm]
\begin{axis}[
x=1.0cm,y=1.0cm,
axis lines=middle,
ymajorgrids=true,
xmajorgrids=true,
xmin=-6.5600000000000005,
xmax=6.640000000000001,
ymin=-3.639999999999998,
ymax=6.440000000000002,
xtick={-6.0,-5.0,...,6.0},
ytick={-3.0,-2.0,...,6.0},]
\clip(-6.56,-3.64) rectangle (6.64,6.44);
\end{axis}
\end{tikzpicture}


\newpage
\def\tocke6{4+4}
\begin{naloga}[\hfill \quad $\tocke6$
\qquad \rightsquigarrow \vert a.\qquad| b.\qquad| c.\qquad| ] \addtocounter{tocke}{\tocke6}
V geometrijskem zaporedju je prvi člen enak 5, četrti člen pa 40. 

a) Izračunaj količnik zaporedja in tretji člen.

b) Izračunaj vsoto prvih 10 členov tega zaporedja.
\end{naloga}

\vspace{10cm}
\def\tocke6{4+3}
\begin{naloga}[\hfill \quad $\tocke6$
\qquad \rightsquigarrow \vert a.\qquad| b.\qquad| c.\qquad| ] \addtocounter{tocke}{\tocke6}
a) Naj bo $\log a=3, \log b=-2,\log c=4.$ Izračunaj $\log x,$ če je $x=\dfrac{a^2b}{\sqrt[4]{c}}.$

b) Izračunaj: $\log_3 9-\log 1-\ln e+\log 2+\log 5$
\end{naloga}
 



\newpage
\def\tocke3{3+4+3}
\begin{naloga}[\hfill \quad $\tocke3$
\qquad \rightsquigarrow \vert a.\qquad\vert b.\qquad] \addtocounter{tocke}{\tocke3}
a) V aritmeričnem zaporedju $-5,\;\rule{1cm}{0.1mm},1,\;\rule{1cm}{0.1mm}\;\rule{1cm}{0.1mm}, \ldots$ dopolni tri manjkajoče člene.

b) Izračunaj 30. člen in vsoto prvih 30 členov.

c) Na katerem mestu je člen $142\;?$
\end{naloga}



\vfill
\mdtheorem[style=theoremstyle1]{kriterij}{Kriterij ocenjevanja}
\begin{kriterij*}[\hfill \v{s}tevilo vseh to\v{c}k na testu: $\arabic{tocke}$]
 $$
 \begin{array}{||c|c|c|c|c|c||
>{}c|c|}
 \hline\hline
 \textup{ocena}& 1&2&3&4&5&\textrm{uspe\v{s}nost v } \%& \textrm{\bf OCENA}\\
  \hline\hline
 \% & [0,45) &[45,60)&[60,75
 ) &[75,90)    &    [90,100]& \mbox{\phantom{ \huge{$\frac{5}{445}$
 15}}
 } 
 & \\
 \hline
 \end{array} \hspace*{0.5cm}\vspace*{-0.3cm} \qrcode[hyperlink,height=0.8in]{http://www.matej.info}
$$
\vspace*{0.1cm}
\end{kriterij*}

\newpage



\subsection{Rešitve:}
\subsubsection{1.}
 a) $\log_3(2x-1)=2$

\[
2x-1=3^2
\]
\[
2x-1=9
\]
\[
2x=10
\]
\[
x=5
\]

Pogoj: ($2x-1>0 \Rightarrow x>\tfrac12$), rešitev je veljavna.

\[
\boxed{x=5}
\]


b) $3^x=6$

Zapišemo z logaritmom:
$$
x=\log_3 6
$$

$$
\boxed{x=\log_3 6}
$$

 c) \sqrt{2}=2^{x-\frac12}

$$
\sqrt{2}=2^{\frac12}
$$

$$
2^{\frac12}=2^{x-\frac12}
$$

Primerjamo eksponente:
$$
\frac12=x-\frac12
$$
$$
x=1
$$

$$
\boxed{x=1}
$$

\subsubsection{2.}

## Graf funkcije $f(x)=3^x-3$

 a) Začetna vrednost, ničla, asimptota

Začetna vrednost:
$$
f(0)=3^0-3=1-3=-2
$$

Ničla:
$$
3^x-3=0
$$
$$
3^x=3
$$
$$
x=1
$$

Vodoravna asimptota:
$$
y=-3
$$

b) Točka $A(2,y)$

$$
y=f(2)=3^2-3=9-3=6
$$

$$
\boxed{A(2,6)}
$$


\begin{tikzpicture}
\begin{axis}[
    axis lines=middle,
    xlabel={$x$},
    ylabel={$y$},
    xmin=-3, xmax=4,
    ymin=-4, ymax=8,
    grid=both,
    width=10cm,
    height=7cm,
    samples=200
]

% funkcija
\addplot[blue, thick] {3^x - 3};

% asimptota
\addplot[dashed, red] {-3};

% točka A
\addplot[only marks, mark=*] coordinates {(2,6)};
\node[right] at (axis cs:2,6) {$A(2,6)$};

% začetna vrednost
\addplot[only marks, mark=*] coordinates {(0,-2)};
\node[left] at (axis cs:0,-2) {$(0,-2)$};

% ničla
\addplot[only marks, mark=*] coordinates {(1,0)};
\node[above] at (axis cs:1,0) {$(1,0)$};

\end{axis}
\end{tikzpicture}
```
\subsubsection{3.}


Podano je:
$$
a_1=5,\qquad a_4=40
$$

a) Količnik in tretji člen

Za geometrijsko zaporedje velja:
$$
a_n=a_1 k^{n-1}
$$

Za četrti člen:

$$
a_4=5k^3=40
$$

$$
k^3=8 \quad\Rightarrow\quad k=2
$$

Tretji člen:
$$
a_3=a_1 k^2=5\cdot 2^2=20
$$

$$
\boxed{k=2,\qquad a_3=20}
$$

 b) Vsota prvih 10 členov

Formula za vsoto prvih $n$ členov:
$$
S_n=a_1\frac{k^n-1}{k-1}
$$

$$
S_{10}=5\cdot\frac{2^{10}-1}{2-1}=5(1024-1)=5\cdot1023=5115
$$

$$
\boxed{S_{10}=5115}
$$

\subsubsection{4.}

a) Izračun $\log x$

Podano:
$$
\log a=3,\quad \log b=-2,\quad \log c=4
$$
$$
x=\frac{a^2 b}{\sqrt[4]{c}}=\frac{a^2 b}{c^{1/4}}
$$

Uporabimo logaritemska pravila:
$$
\log x=\log a^2+\log b-\log c^{1/4}
$$
$$
\log x=2\log a+\log b-\frac14\log c
$$
$$
\log x=2\cdot3+(-2)-\frac14\cdot4
$$
$$
\log x=6-2-1=3
$$

$$
\boxed{\log x=3}
$$

b) Izračun izraza

$$
\log_3 9-\log 1-\ln e+\log 2+\log 5
$$

$$
\log_3 9=2,\quad \log 1=0,\quad \ln e=1
$$

$$
\log 2+\log 5=\log(10)=1
$$

$$
2-0-1+1=2
$$

$$
\boxed{2}
$$

\subsubsection{5.}

Podano zaporedje:
$$
-5,\;\rule{1cm}{0.1mm},\;1,\;\rule{1cm}{0.1mm},\;\rule{1cm}{0.1mm},\ldots
$$

### a) Manjkajoči členi

Razlika:
$$
d=\frac{1-(-5)}{2}=3
$$

Zaporedje:
$$
-5,,-2,,1,,4,,7,\ldots
$$

b) $30$. člen in vsota prvih $30$ členov

$$
a_n=a_1+(n-1)d
$$

$$
a_{30}=-5+29\cdot3=82
$$

Vsota:
$$
S_{30}=\frac{30}{2}(a_1+a_{30})=15( -5+82)=15\cdot77=1155
$$

$$
\boxed{a_{30}=82,\quad S_{30}=1155}
$$

c) Mesto člena 142

$$
142=-5+(n-1)3
$$
$$
147=3(n-1)
$$
$$
n-1=49 \Rightarrow n=50
$$

$$
\boxed{142 \text{ je } 50.\text{ člen}}
$$

neobvezno: graf
\begin{tikzpicture}
\begin{axis}[
    axis lines=middle,
    xlabel={$n$},
    ylabel={$a_n$},
    xmin=0, xmax=11,
    ymin=-6, ymax=25,
    grid=both,
    width=10cm,
    height=7cm
]
\addplot[
    only marks,
    mark=*,
    blue
]
coordinates {
(1,-5) (2,-2) (3,1) (4,4) (5,7)
(6,10) (7,13) (8,16) (9,19) (10,22)
};
\end{axis}
\end{tikzpicture}



\end{document}